\begin{corollary}[6.15.9]
\label{IV.2.s6.15.9-fr}
Sean $k$ un cuerpo, $X$ un $k$-preesquema localmente de
tipo finito sobre $k$ y $k'$ una extensión separablemente
cerrada de $k$ \textup{(es decir, tal que la clausura algebraica de
$k'$ sea radicial sobre $k'$)}. Entonces, $X$ es
geométricamente unirramoso si y solo si
$X'=X\otimes_k k'$ es unirramoso.
\end{corollary}

\begin{proof}
En vista de (6.15.7), basta probar que, si $X$ es unirramoso
y $k$ es separablemente cerrado, entonces $X$ es
geométricamente unirramoso. Ahora bien, $X$ es
geométricamente unirramoso en sus puntos cerrados (6.15.8);
deduciremos en (10.4.9) que $X$ es geométricamente unirramoso
en todos sus puntos, apoyándonos en el hecho (señalado al
final de (6.15.2)) de que el conjunto de los puntos donde
$X$ es geométricamente unirramoso es construible (por
supuesto, el corolario (6.15.9) no se utilizará hasta
entonces).
\end{proof}

Este resultado justifica, en cierta medida, la terminología
«geométricamente unirramoso».
